\input{../tex-inputs/latex-leseansicht-vorspann}

\section[Emil Marriot an Arthur Schnitzler, 9. 5. 1902]{L04335 Emil Marriot an Arthur Schnitzler, 9. 5. 1902}
\nopagebreak\mylabel{L04335v}
\rehead{ }\normalsize\beginnumbering\briefempfaengerindex{Schnitzler, Arthur@\textsc{Schnitzler, Arthur}!zzzMarriot, Emil@\emph{von Emil Marriot}!1902-05-091@{9. 5. 1902}|(be}
\toendnotes[C]{\smallbreak\pagebreak[2]}
\correspDesc{Versand  durch Emil Marriot am 9. 5. 1902 in Wien
\newline{}Erhalt  durch Arthur Schnitzler im Zeitraum [10. 5. 1902
                  – 14. 5. 1902?] in Wien}\toendnotes[C]{\smallbreak}
\Standort{DLA, A:Schnitzler, HS.1985.1.4008.}
\physDesc{Brief, 1 Blatt, 2 Seiten, 762 Zeichen
\newline{}Handschrift: schwarze Tinte, lateinische Kurrent}\toendnotes[C]{\smallbreak}
\pstart
           \raggedleft{}{\pb}Wien\oindex{Wien@\textbf{Wien}, \emph{Verwaltungsgebiet}|pw}, 9. V. 1902.\pend
           
\pstart{}Sehr gehrter Herr,\pend\vspace{0.5em}
\pstart
           passt Ihnen der nächste \label{K_L04335-1v}\edtext{Dinstag}{\lemma{\textnormal{\emph{Dinstag}}}\Cendnote{\textnormal{Vgl. A. S.: \emph{Tagebuch}, 13. 5. 1902.}}}\label{K_L04335-1}? So
               zwischen ½ 12 – 12 Uhr? Natürlich steht es Ihnen völlig frei, einen andern Vormittag
               der nächsten Woche zu wählen. Ich habe den Dinstag nur ganz zufällig ins Auge
               gefasst, weil ich irgend einen besti{\geminationm}ten Tag nennen
               will. Nur bitte ich Sie, mich vorher wissen zu lassen, an welchem Tag ich die Freude
               haben werde, Sie zu begrüssen. Wenn keine Veranlassung vorliegt, {\pb}zu Hause zu bleiben, gehe ich um diese Zeit gewöhnlich
               aus. Der Arbeit gehört auch bei mir der Nachmittag. Aber derzeit bin ich unfähig, zu
               schreiben.\pend
           
\pstart
           Entschuldigen Sie, bitte, dass ich in Allem so rasch bin. Ich bleibe nicht mehr lang
               hier, höchstens bis \label{K_L04335-5v}\edtext{AnfangJuni}{\lemma{\textnormal{\emph{AnfangJuni}}}\Cendnote{\textnormal{Sie reiste auf das Schloss Brodzany\oindex{XXXX Ortsangabe fehlt|pwk} (Brogyán\oindex{XXXX Ortsangabe fehlt|pwk}), vgl. XXXX Auszeichnungsfehler: Dokument L04336 nicht gefunden.}}}\label{K_L04335-5}. Darum die Eile.\pend
           
\pstart
           Mit verbindlichen Grüssen{\\[\baselineskip]}Ihre ergebene{\\[\baselineskip]}\spacefill\mbox{E. Marriot}\pend
           \leftskip=0em{}\selectlanguage{ngerman}\endnumbering\briefempfaengerindex{Schnitzler, Arthur@\textsc{Schnitzler, Arthur}!zzzMarriot, Emil@\emph{von Emil Marriot}!1902-05-091@{9. 5. 1902}|)be}\mylabel{L04335h}
\begin{anhang}
\end{anhang}\newcommand{\dateiname}{L04335}\newcommand{\titel}{Emil Marriot an Arthur Schnitzler, 9. 5. 1902}\newcommand{\editorInnen}{Herausgegeben von Jahnke, SelmaMüller, Martin Anton}\input{../tex-inputs/latex-leseansicht-abspann}
